% Figure: Mohr's circle for the worked stress state sigma_xx=80, sigma_yy=40,
% sigma_xy=30. Center at (60,0), radius = sqrt(20^2 + 30^2) = 36.06.
% sigma_1 = 96.06, sigma_2 = 23.94. Points X=(80,30), Y=(40,-30).
\begin{figure}[htbp]
\centering
\begin{tikzpicture}
\begin{axis}[
  width=12cm, height=9cm,
  axis equal image,
  xlabel={normal stress $\sigma$ (MPa)},
  ylabel={shear stress $\tau$ (MPa)},
  xmin=10, xmax=110, ymin=-50, ymax=50,
  axis lines=middle,
  tick label style={font=\scriptsize},
  label style={font=\small},
]
% circle: center (60,0), radius 36.06
\draw[engblue, very thick] (axis cs:60,0) circle [radius=36.06];
% center
\filldraw[enggray] (axis cs:60,0) circle [radius=1.2pt];
\node[enggray, font=\scriptsize, anchor=north] at (axis cs:60,-2) {$C(60,0)$};
% point X = (80,30)
\filldraw[engred] (axis cs:80,30) circle [radius=1.6pt];
\node[engred, font=\scriptsize, anchor=south west] at (axis cs:80,30) {$X(80,30)$};
% point Y = (40,-30)
\filldraw[engred] (axis cs:40,-30) circle [radius=1.6pt];
\node[engred, font=\scriptsize, anchor=north east] at (axis cs:40,-30) {$Y(40,-30)$};
% diameter XY
\draw[engred, thick] (axis cs:80,30) -- (axis cs:40,-30);
% principal stresses on axis
\filldraw[engorange] (axis cs:96.06,0) circle [radius=1.6pt];
\node[engorange, font=\scriptsize, anchor=south] at (axis cs:96.06,3) {$\sigma_1=96.1$};
\filldraw[engorange] (axis cs:23.94,0) circle [radius=1.6pt];
\node[engorange, font=\scriptsize, anchor=south] at (axis cs:23.94,3) {$\sigma_2=23.9$};
% tau max at top
\filldraw[englime!60!black] (axis cs:60,36.06) circle [radius=1.6pt];
\node[font=\scriptsize, anchor=south] at (axis cs:60,38) {$\tau_{\max}=36.1$};
\end{axis}
\end{tikzpicture}
\caption[Mohr's circle for the worked stress state]{Mohr's circle for
$\sigma_{xx}=80$, $\sigma_{yy}=40$, $\sigma_{xy}=30$ MPa. The $x$-face
plots as $X=(\sigma_{xx},\sigma_{xy})=(80,30)$ and the $y$-face as
$Y=(\sigma_{yy},-\sigma_{xy})=(40,-30)$; the two are diametrically
opposite. The circle has centre $C=((\sigma_{xx}+\sigma_{yy})/2,0)=(60,0)$
and radius $R=\sqrt{20^2+30^2}=36.1$ MPa. The intersections with the
horizontal axis are the principal stresses $\sigma_1=96.1$ and
$\sigma_2=23.9$ MPa; the top of the circle is the maximum in-plane shear
$\tau_{\max}=R=36.1$ MPa. Rotating the physical axes by $\theta$ rotates
the diameter $XY$ by $2\theta$ on the circle.}
\label{fig:vol03ch03:mohr-circle}
\end{figure}
